\section{Hall reduction: absolute densities and the mixed angular law}
\label{sec:hall}

A deficient matching leaves macroscopic colors at sites that cannot be
treated by the low-deficit angular expansion. We first identify the positive
measure carried by these sites through an absolute determinant calculation.
We then prove that the finite, signed angular transfer joins it to the
dormant matching law. Throughout this section \(r=k+2\),
\(Z_r=S_{11}\tau_r\), and the Lie algebra is in unit-quaternion coordinates,
with inner product \(-\operatorname{Tr}(XY)/2\). In particular,
\(\operatorname{Vol}(\mathrm{SU}(2))=2\pi^2\) and the unit axial circle has
period \(2\pi\). All seam signs and fibre periods refer to the stated
untwisted paired-cap surgery.

\subsection{Intrinsic Hall data and the active cap measure}
\label{sec:hall:data}

Let \(N\) be the number of backbone vertices, \(\nu\) its matching number,
and let \(S_*\) consist of the vertices exposable by some maximum matching.
Set
\[
 T_*=N_G(S_*),\qquad I=G\setminus(S_*\cup T_*),\qquad
 a=|T_*|,\quad b=|S_*|,\quad d=|I|/2,\qquad \nu=a+d .
\]
The inactive forest \(I\) has a unique perfect matching \(M_I\).
Equivalently, begin with a maximum matching saturating a minimum cover
and follow alternating paths from the unmatched vertices in its
complement. The reached complement vertices are precisely \(S_*\), and
their reached cover neighbors are \(T_*\). Each \(t\in T_*\) has at least
two distinct neighbors in \(S_*\): its matching partner and its
predecessor on a path from an exposed vertex. Also, no edge joins \(S_*\)
to \(I\). These properties give the minimum cover \(T_*\cup A_I\), where
\(A_I,B_I\) are the two bipartitions of the dormant forest. The
bipartitions may be chosen independently in each component. If \(a_jb_j\)
is the \(j\)-th matched edge, an unmatched \(a_i b_j\) gives the relation
\(i\to j\). Write
\[
 \Omega_I=\{u\in[0,1]^d:u_i\le u_j\text{ whenever }i\to j\},\qquad
 \iota_I(u)=(u_j,1-u_j)_{j=1}^d .
\]
The coordinates in \(\iota_I\) retain their physical vertex labels.
Reversing a component's bipartition reflects its root coordinates and
does not change any resulting labelled measure.

For a noncentral fibre \(h_s\) put \(x_s=1-\arg(h_s)/\pi\).
At \(t\in T_*\), denote by \(P_t\) the ordered product of the incoming
noncentral conjugacy-class variables, with the prescribed central
factors included. Set
\begin{equation}
 dM_*=\prod_{s\in S_*}2\cos^2(\pi x_s/2)\,dx_s\prod_{t\in T_*} L_t(x),
 \qquad
 L_t(x)=\mathbb E\!\left[\frac{\pi}{4|\operatorname{Im}P_t|}\right].
 \label{eq:hall:weight}
\end{equation}
The expectation uses independent probability conjugacy-class variables at
each cut active block. Successive tree gluing fixes the relative
orientations; it does not replace \(P_t\) by the product of its traces.

Here is the accompanying angular coarea rule. For the two trace-zero
cap quaternions \(A=u,B=v\) with probability sphere law \(d\Pi\), the
real part of \(Q=AB\) is \(-u\cdot v\), uniform on \([-1,1]\).
The group-Haar density of \(Q\) is therefore
\(\pi/(4|\operatorname{Im}Q|)\). At \(Q\ne\pm1\), rotation about its
axis gives the conditional uniform law of \(u\) on its perpendicular
equator. In a common local frame the resulting cap kernel at \(P_t\)
can be written
\begin{equation}
              \frac{\pi}{2}\,
          \delta\bigl(u_t\cdot\operatorname{Im}P_t\bigr)\,d\Pi(u_t).
 \label{eq:hall:cap-coarea}
\end{equation}
In fact the integral of \(\delta(u\cdot q)\) against \(d\Pi(u)\) is
\(1/(2|q|)\) for \(q\ne0\). At a central product this formula is used as
coarea in the variables of the ordered group relation, not as a
normalized pointwise equator. The conditional circle generally depends
on the direction from which the product approaches the center.

The same \(L_t\) is forced by the finite quantum characters. With
\(\chi_n(h)=\sin(n\arg h)/\sin\arg h\), the coefficient of \(\chi_n\) in
\(\pi/(4\sin\arg h)\) under probability group Haar is \(1/n\) for odd \(n\)
and zero for even \(n\). Averaging a product character over independent
conjugacy-class representatives gives
\(n\prod_s(\chi_n(h_s)/n)\), where the product contains one factor
\(\chi_n(h_s)/n\) per incoming active vertex. Thus on compact interior
macro radii
\begin{equation}
 L_t(x)=\sum_{\substack{n\ge1\\ n\ {\rm odd}}}
             \prod_{s\in S_*\cap N(t)}\frac{\chi_n(h_s)}n .
 \label{eq:hall:characters}
\end{equation}
The convergence is absolute there because \(t\) has at least two incoming
active neighbors. Formula \eqref{eq:hall:cap-coarea}, rather than the
pointwise endpoint values of \eqref{eq:hall:characters}, defines the
positive measure used in endpoint exhaustion.

\subsection{The absolute adjoint determinant line}
\label{sec:hall:torsion}

We give the normalization on a slightly more general chart, in order to
state the mixed-sign overlap without specializing the number of sites.
Let \(S\) be an independent noncentral set, \(T=N_G(S)\) and
\(I_0=V\setminus(S\cup T)\), of sizes \(b_0,a_0,n_0\).
Assume each \(t\in T\) is negative central and its actual cap tuple is
locally irreducible; this holds in particular for noncentral \(P_t\).
Work first where its fixed-periphery volume is regular, and on globally
irreducible characters. Let \(n_+\) be the number of positive centers in
\(I_0\), fix their cap-root choices, and sum the two active roots at
each \(s\). At each inactive site its invariant coefficient space \(V_i\)
is the line of its negative cap axis or the full quaternion algebra at a
positive center. Let \(\mathcal B_{I_0}\) be the symmetric projected
adjacency on \(\bigoplus_{i\in I_0}V_i\) given by the marked seam maps,
and set \(K=\ker\mathcal B_{I_0}\). All these spaces carry the indicated
unit metric. Use \(dM_{S,T}\) for the analogue of \eqref{eq:hall:weight}
and \(dP\) for the conditional normalized active-surface and bending
Haar variables. The latter are genuine cut variables on the effective
group quotient. In cohomology, use the dual \(H^2\) basis to the actual
deformation basis of \(H^1\).

\begin{theorem}[Absolute Hall determinant factorization]
\label{thm:hall:torsion}
The stack half-torsion in these coordinates is
\begin{equation}
 d\mu_{\rm stack}
 =2^{a_0+N-1-4n_+}\pi^{\,2b_0+n_0-2-n_+}\,
    dM_{S,T}\,dP\,
    |\det{}'\mathcal B_{I_0}|^{-1/2}\,d\operatorname{vol}_{K}.
 \label{eq:hall:torsion}
\end{equation}
Here \(\det{}'\) is the product of nonzero eigenvalues, with empty
product one, and \(d\operatorname{vol}_{K}\) uses the actual unit
quaternion fibre periods. The formula gives the density in all retained
harmonic variables on a fixed-rank chart. It does not claim that each
harmonic vector at a singular reference is tangent to a flat family.
If the positive roots are also summed, multiply by \(2^{n_+}\).
\end{theorem}

\begin{proof}
Use ordinary unit coefficient cellular cochains, in which a differential
from degree \(j\) to \(j+1\) contributes the absolute determinant to
power \((-1)^j\). A torus with a noncentral peripheral holonomy has only
its invariant line in cohomology, in primitive section and fibre units;
the perpendicular complex is acyclic with torsion one. A central seam
has the same unit assertion in all three directions.

At a noncentral vertex of valence \(k\), the cap-free piece is a
punctured sphere times the fibre circle. The perpendicular fibre
differential has determinant \(4\sin^2(\pi x)\), and its Euler exponent
from the base is \(-k\). Each cap solid torus contributes a
perpendicular determinant \(4\cos^2(\pi x/2)\). On the axis the product
cochains have canceling degrees. At each filling the section difference
has unit pivot and the primitive fibre evaluates by \(\pm2\) on the
core. The degree-one cap mismatch therefore gives \(1/2\) in torsion
per cap; the degree-zero and degree-two pivots are one. Thus the piece
half-torsion is
\[
                  \frac{2\cos^2(\pi x/2)}{(2\sin(\pi x))^k}.
\]
At a negative central piece the two cap cores act by perpendicular
half-turns, each with determinant \(4\). The same axial primitive
subdeterminants give \(4/2\) per cap in torsion, hence the piece
half-torsion is \(2\). In the inactive positive case the coefficient
system is trivial and the axial cap calculation occurs in three
directions, giving half-torsion \(1/8\).
For a locally irreducible negative piece, choose the cap-free base
cohomology in paired degrees using all peripheral traces, possible by
relative duality since the invariants vanish. The cap determinants
then remove exactly their own two trace lines. The same formula \(2\)
holds, with fixed-periphery cohomology retained.

The remaining Mayer--Vietoris calculation can be made with primitive
half-edge section coordinates. In degree one augment each vertex
restriction by its one section sum in the local invariant space.
Eliminate the section mismatches with identity pivots. A noncentral
vertex keeps its one common fibre period; its incident section values
use a sum-zero basis such as \(e_i-e_k\), with unit complementary lift
\(e_k\). A locally irreducible active center has no invariant fibre
period. Since \(S\) is independent and no \(S\)--\(I_0\) edge exists, the
remaining fibre adjacency is zero on the old \(S\) periods and is
exactly \(\mathcal B_{I_0}\) on the inactive periods. In degree zero the
map is the inclusion-difference \(U\) of vertex invariant lines or
algebras in the transported seam spaces. The degree-two half-edge
augmentation gives \(U^*\): one-sided lifts of the seam area sums have
determinant one, so no norm of a diagonal or sum-zero lattice is
inserted. The connecting collar cup product pairs each retained unit
fibre period with its own unit quotient dual, and the fixed-periphery
pairs use the Goldman form. Cross terms from alternate lifts are
triangular, since the connecting collars pair through the indicated
boundary periods. In the unit fibre metric the residual determinant is
therefore \(|\det{}'\mathcal B_{I_0}|^{-1/2}\), with its induced kernel
volume. This specifies the volume on a slanted kernel rather than
replacing it by the volume of any selected coordinate coefficients.

For completeness, put \(p=\sum_{s\in S}\deg s\) and \(q=N-1-p\). The
\(p\) noncentral seam circles and \(q\) central seam groups contribute
\((2\pi)^p(2\pi^2)^q\). Their effective quotient follows from the paired
\(U,U^*\) determinants, or by inserting the vertex invariant generators
as columns and taking the determinant complement. Divide by the
\(b_0\) noncentral vertex circles,
the \(n_0-n_+\) inactive circles and \(n_+\) inactive groups, and the
\(a_0\) locally irreducible centers.
The resulting stack angular factor is
\[
 \frac{2^{-a_0}(2\pi)^{p-b_0}(2\pi^2)^q}
      {(2\pi)^{n_0-n_+}(2\pi^2)^{n_+}}.
\]
Indeed the kernel of the combined local conjugation is precisely the
single global center when the closed tuple is irreducible. Dividing by
the product of the full local stabilizers includes that final stack
factor; dividing by it a second time would be incorrect.

The active unit period changes to \(\pi\,dx_s\). If \(k_t\) is the
number of active neighbors of a negative center, its absolute reduced
Goldman volume is
\[
   2^{k_t+1}\pi^{k_t-2}
        \left(\prod_{s\sim t,\ s\in S}\sin(\pi x_s)\right)L_t .
\]
For this last normalization use the group-valued volume formula
\cite{AMW}: a noncentral conjugacy class has volume
\(2\pi\sin\theta\); the two traceless classes each have volume \(2\pi\);
and the coarse fixed-product reduction uses
\(2/\operatorname{Vol}(\mathrm{SU}(2))^2\) times the convolution at
identity in these probability-Haar conventions. Formula
\eqref{eq:hall:cap-coarea} supplies its convolution factor.
Equivalently this is the relative cup pairing on the punctured
sphere followed by one primitive fibre product. Multiplying these
volumes and all block factors, cancel the active incidence sines,
use \(\sum_t k_t=p\) and \(p+q=N-1\), and include the \(2^{b_0}\)
true noncentral root sheets. This gives \eqref{eq:hall:torsion}.
\end{proof}

In particular let \(S=S_*,T=T_*,I_0=I\), with all dormant signs
negative. For each dormant edge let \(t_e\) be the actual signed cap
axis dot. The determinant of the weighted bipartite adjacency of a
perfectly matched forest satisfies
\[
          |\det\mathcal B_I|^{1/2}=\prod_{e\in M_I}|t_e| .
\]
Every nonzero determinant monomial gives a perfect matching and two
different matchings on a forest would have an alternating cycle.
On its resolved regular charts this gives
\begin{equation}
 d\mu_{\rm reg}
      =C_{\rm abs}\,dM_*\,\frac{dP}{\prod_{e\in M_I}|t_e|},
 \qquad C_{\rm abs}=2^{a+N-1}\pi^{\,2b+2d-2}.
 \label{eq:hall:resolved-density}
\end{equation}
When a dormant matched component is orthogonal and both radial
periods at every matched edge are retained, the block determinant is
instead the corresponding literal \(\pi\,dy_i\) per signed fibre
defect, using exactly the same angular density.

Rooting the central seams makes the conditional joint dormant axes
transparent. The incoming central relative rotation of a new dormant
axis has independent group Haar law and so produces a uniform sphere
axis conditional on earlier axes, local active cut shapes and active
radii. The old ordered products see their dormant neighbors' central
fibre signs, not these cap axes. Thus the pre-Fourier dormant axis
measure is the product of spherical probabilities, modulo the common
rotation and actual cap sign actions. Its dots on tree edges are
independent uniform variables. Conditioning on \(t_e=0\) for each
\(e\in M_I\) gives an ordered orthogonal frame on that edge and the
joint law of independent Haar frames. A joint Wilson word is evaluated
in these shared axes, not in separate angle marginals. This description
does not impose independence conditional on arbitrarily prescribed
global relative orientations of old active blocks.

For the actual local action and its calibrated stable reduction
through arbitrary dormant angles and interior active locks, see
Theorems~\ref{classic:thm:hall-action} and
\ref{classic:thm:hall-stationary}. In particular the \(2\pi^2\)
conversion between the quaternion-unit curl and the action Hessian
compensates the radial coarea factor \(\pi^{2d}\). Applying the
Gaussian normalization to \eqref{eq:hall:resolved-density} gives the
absolute factor that will occur in the joint quantum law:
\begin{equation}
 p_G=(3N-2\nu-3)/2,\qquad
 c_G=C_{\rm abs}(2\pi^2)^{-p_G}
    =2^{(1-N)/2+2a+d}\pi^{1-N}.
 \label{eq:hall:normalization}
\end{equation}
The signed phase follows independently from the capped flat action;
on this Hall stratum it is \((-1)^{r(N-1)}\).

\subsection{Active endpoints and bounded-product overlap}
\label{sec:hall:overlap}

The cap measure has a useful uniform bound at genuine macro
endpoints. Put \(\alpha_s=\min(x_s,1-x_s)\). On the open cube,
\[
       L_t(x)\asymp_G
        \frac{1}{\max_{s\in S_*\cap N(t)}\alpha_s}.
\]
For the upper bound average last the class with largest reduced
angle \(\alpha\), using radians for the following calculation. For a
fixed other quaternion of angle \(\beta\), the angle of their product
has density \(\sin\theta/(2\sin\alpha\sin\beta)\) on an interval of
length \(2\min(\bar\alpha,\bar\beta)\), where a bar reduces angles to
\([0,\pi/2]\) by central signs. This follows directly from the uniform
relative axis dot, and extends to a central fixed input by
continuity. The integral of \(1/\sin\theta\) is bounded by a constant
times \(1/\sin\alpha\). The lower bound uses the triangle inequality
on the distance of the product to the center.

For \(\alpha_s\asymp2^{-k_s}\), let \(A_j=\{s:k_s\ge j\}\) and
\(B(A_j)=\{t:N_{S_*}(t)\subset A_j\}\). The mass of the dyadic box in
\eqref{eq:hall:weight} is bounded by a constant times
\[
 2^{-\sum_s k_s+\sum_t\min_{s\sim t}k_s}
       =2^{-\sum_{j\ge1}(|A_j|-|B(A_j)|)} .
\]
For nonempty \(A_j\), the forest on \(A_j\) and \(B(A_j)\) has at least
\(2|B(A_j)|\) incidences and at most \(|A_j|+|B(A_j)|-1\).
Thus the exponent is at least \(\max k_s\). Summing by that maximum
gives, for \(b>0\),
\begin{equation}
 \int_{\{\min_s\alpha_s<\epsilon\}}dM_*
 \le C_G\epsilon(1+|\log\epsilon|)^{b-1}.
 \label{eq:hall:classical-end}
\end{equation}
The same absolute bound holds with any bounded joint classical angular
mark. There is also no endpoint mass hidden at a lock of \(P_t\).
On \(\alpha_s\ge\epsilon\) the product of two incoming classes has
group-Haar density bounded by \(C\epsilon^{-2}/\sin\theta\).
The remaining input products translate this bound. The integral of
the translated inverse-distance factor times \(1/\sin\theta\) over a
ball of radius \(\delta\) at either center is \(O(\epsilon^{-2}\delta)\),
uniformly over the translate. If the two poles are close to the ball,
apply Cauchy--Schwarz and the \(O(\delta)\) inverse-square integrals in
dimension three; otherwise integrate directly. Bounding the other
\(a-1\) cap factors by \(C\epsilon^{-(a-1)}\), and using
\eqref{eq:hall:classical-end} on the complement, gives a uniform
vanishing mass, for example
\(O(\delta^{1/(a+2)}(1+|\log\delta|)^{b-1})\), for the union of these
locks when \(a>0\). The statement does not require a pointwise choice
of conditional angular law at a collapsed equator.

We next make the normalization of all bounded physical transition
regions explicit. Let \(W\) be a fixed collection of integer-spin
Wilson words piercing only positive dormant cap disks. The words may
join dormant components. Use the exterior-compatible tree-corridor
routing: auxiliary tracks run on one side of each clasp corridor,
above the old forced sublink, and descend only at their dormant cap
ports. In the diagram with the external \(T_*\)--\(I\) Hopf clasps
removed, the ordinary forced sublink can be placed below a plane
separating it from the dormant and auxiliary ribbons. This is a
specified class of actual framed links, not an identification of all
embeddings of an abstract word. It satisfies the cut hypotheses of
Lemma~\ref{lem:ang:network}. The cap-word evaluation on the central
dormant axes is \(P_W\) in \eqref{eq:ang:angular-kernel}. In particular,
there are no old active cap passages or deliberately retained
linkings with a macroscopic ribbon in the separated diagram.

There is also an exact stronger cut identity for words routed
entirely inside individual connected dormant components. Contract
those components to vertices of the incidence tree and prune leaves.
Cut the one simple backbone ribbon at the exterior clasp, keeping
the clasp outside the cut. A closed decorated leaf then gives an
endomorphism of a simple object, hence its normalized standalone
scalar. Reattaching contributes just the ordinary normalized Hopf
factor. Repeating the cut gives exactly the product of the standalone
decorated dormant transfers and the unchanged outside scalar Hopf
factors, at every level. This argument never divides by a scalar
Hopf entry, which can vanish. A connecting word can place a nonscalar
auxiliary packet on that cut, which is why its proof uses the
conditional forced exhaustion below.

Let \(E'\) be any subset of the unmatched dormant arrows. Keep a fixed
compact positive product annulus with smooth profile \(\psi_j(w_j)\)
at each physical matched scale
\(w_j=m_{a_j}m_{b_j}/r\), and a fixed compact positive profile
\(\chi_e(v_e)\) at each \(v_e=m_{a_i}m_{b_j}/r\), \(e=(i\to j)\in E'\).
Contract the \(E'\) arrows on the matching tree. Let \(p=d-|E'|\),
let \(\Omega_{E'}\) be the order polytope on the \(p\) components, and
let \(\iota_{E'}\) insert the corresponding exponent pair
\((u_C,1-u_C)\) at every physical matched pair in its component $C$.
At the other unmatched edges one may use fixed bounded upper-product
gates equal to one near zero; alternatively support the smooth
logarithmic mark \(F\) in the strict interior chamber for the
remaining order comparisons. Use fixed negative-centre profiles
equal to one near zero on the forced and dormant sites and a fixed
smooth macro mark \(\eta\) with compact interior support.

\begin{theorem}[Absolute mixed overlap at bounded physical products]
\label{thm:hall:overlap}
For the compatible literal words and these marks,
\begin{align}
 &\lim_{r\to\infty}
 \frac{(-1)^{r(N-1)}Z_{r,W}[\eta,F,\psi,\chi]}
      {2^d c_G\,r^{p_G}(\log r)^p}
 =\left(\int\eta\,dM_*\right)
  \left(\int_{\Omega_{E'}}F(\iota_{E'}(u))\,d^pu\right)
 \notag\\[-0.2em]
 &\quad{}\times
  \int_{\mathbb R_+^{d+|E'|}}
    \mathcal K_{W,I}(\mathbf w_M,\mathbf v_{E'},0_{\rm rest})
    \prod_j\psi_j(w_j)\,dw_j
    \prod_{e\in E'}\chi_e(v_e)\,\frac{dv_e}{v_e}.
 \label{eq:hall:overlap}
\end{align}
Here \(\mathcal K_{W,I}\) is the actual joint kernel of
Theorem~\ref{thm:ang:compact}, on the dormant sphere variables with the
original shared words; it is not a product of single-edge observables.
\end{theorem}

The exhaustion proved below justifies first fixing a finite forced tuple,
also for connecting words. To see the normalization, take
\(A=T_*\cup A_I,\ B=S_*\cup B_I,\ \alpha=|B|=N-\nu\).
If \(\mathcal S_{r,W}\) denotes the sum with vertex factors
\(1+\cos(\pi m_v/r)\) and the normalized categorical kernel of
\eqref{eq:ang:normalized-sum}, put
\(X_{r,W}=(2/r)^\alpha 2^{-\nu}\mathcal S_{r,W}\).
The exact external factor in this normalization is asymptotic to
\[
 (-1)^{r(N-1)}2^{(3N+1)/2-2\alpha}\pi^{1-N}r^{p_G}
   =(-1)^{r(N-1)}2^d c_G r^{p_G}.
\]
At fixed forced labels the macro summations have ordinary odd-grid
Riemann limits with conditional density
\(2\cos^2(\pi x_s/2)\,dx_s\) and one normalized character for every old
active incidence. Summing the forced characters gives
\eqref{eq:hall:characters}. At every dormant matched partner, the
odd grid of \(w_j\) has spacing \(2m_{a_j}/r\), and the normalized
conditional mass is \(2\,dw_j/m_{a_j}\). The annuli on \(E'\) bound
all ratios between anchors in each selected component. Keep its
one root anchor and its remaining edge log ratios. Each nonroot
anchor contributes \(dR_e/(2R_e)\), \(R_e=m_{a_i}/m_{a_j}\);
the log-ratio incidence Jacobian is one, by successive leaf
elimination on a tree. Each root yields \((\log r)\,du_C/2\).
The replacement \(v_e=R_e w_j\) then gives the physical
\(dv_e/v_e\), including when several arrows use the same matched
partner. Since \(p+|E'|=d\), the powers of two cancel exactly.

On separated compact log cells all these meshes tend to zero and
all dormant products in the kernel are bounded. The true packet
limit in Theorem~\ref{thm:ang:compact} is uniform there. On the
allowed order chamber the other products tend to zero; on violated
cells they are removed by the stipulated gates. The boundary strips
have arbitrarily small normalized harmonic mass, with uniformly
bounded compact-annulus conditional transfers. This gives
\eqref{eq:hall:overlap}, first for separated smooth tests and then
for smooth joint masks by the usual tensor expansion or mixed
variation bounds. A single measured intersection of several
positive annuli is integrated in its complete \(E'\) chart. Its
contribution must not be replaced by the sum of independently
integrated single-seam residues.

\subsection{Ungated active-Hall angular gluing}
\label{sec:hall:leading}

\begin{theorem}[Ungated literal joint moments for the intrinsic Hall sector]
\label{thm:hall:main}
Let \(\eta\) be any fixed smooth joint test of the literal ordinary
radii of \(S_*\), including both endpoints, and let \(F\) be any
fixed smooth joint test of the \(2d\) literal logarithmic radii of
\(I\). For the exterior-compatible literal words above,
\begin{align}
 Z_{r,W}[\eta,F]
 &=(-1)^{r(N-1)}c_G\,r^{p_G}(\log r)^d
 \left[
       \left(\int\eta\,dM_*\right)
       \left(\int_{\Omega_I}F(\iota_I(u))\,du\right)
       \mathbb E_{\rm frame} P_W+o(1)\right].
 \label{eq:hall:main}
\end{align}
There are no original forced-color prefixes, radial endpoint,
product-scale, or Abel gates. The expectation is the genuine
joint law of independent ordered orthogonal Haar frames on
\(M_I\), modulo common rotation and the cap sign actions.
Ordinary smooth radius arguments on \(T_*\cup I\), if included in a
joint test, may be evaluated at their all-negative centres in the
limiting coefficient. The theorem concerns fixed tests of the
signed invariant, not convergence in unsigned total variation.
\end{theorem}

\begin{proof}
Use the finite-rank network of Lemma~\ref{lem:ang:network} and the
signed full-alcove prefix estimate of Lemma~\ref{lem:ang:prefix}.
Expand in its finitely many rank labels. Fold each color independently
into \([1,r/2]\); the reflected parity is fixed on each branch.
The vertex factors have bounded sup and sampled variation, including
any interval prefix, and the unaltered scalar cap profile is
\(1\pm\cos(\pi n/r)\). Sum \(B\) conditionally first. In \(X_{r,W}\)
normalization a \(B\)-site \(z\) contributes at most \(C/M_z\), where
\(M_z\) is its largest folded neighbor. For its matched anchor
\(n_j\) extract \(1/n_j\), leaving at most \(C n_j/M_z\le C\).
All \(A\)--\(A\) coefficients are bounded. This is an estimate after
each signed \(B\) transfer, not a claim about the absolute variation
of the original color array.

For every \(t\in T_*\) choose the alternating path towards an exposed
vertex. Along that path a matched partner contributes at most a
constant times \(n_{\rm next}/n_{\rm previous}\), and the final
unmatched partner gives \(1/n_{\rm last}\). Their product is bounded
by \(C/n_t\). The factors along each chosen path are distinct, and all
of them lie in the old forced part. Take the geometric mean over
\(T_*\) and restore the exterior harmonic factors. For \(a>0\) the
resulting bound on the remaining \(A\)-summand is
\begin{equation}
 C\prod_{t\in T_*} n_t^{-1-1/a}\,
            \prod_{j\in A_I}n_j^{-1}.
 \label{eq:hall:uniform-forcing}
\end{equation}
If \(a=0\) only the dormant harmonic part is present. This holds on
every folded parity and for all rectangular prefixes. In
particular a forced folded prefix \(n_t\le H\) may be inserted at
uniform normalized cost \(O(H^{-1/a})\). An originally reflected
forced vertex within a bounded folded prefix carries its untouched
\(1-\cos(\pi n_t/r)=O_H(r^{-2})\) factor and is negligible. Thus the
relevant fixed tuple can be taken in the original low-deficit
branch. Mixed variation expansion against prefixes gives the same
statements for the fixed smooth joint tests in the theorem.

This estimate also controls old macro endpoints without a compact
interior assumption. With the other forced labels at most \(H\),
the original endpoint condition at any \(s\in S_*\) restricts its
independent \(B\)-sum to at most \(O(\epsilon r+1)\) entries. Its
normalized transfer is then bounded by
\(C(\epsilon+r^{-1})\le CH(\epsilon+r^{-1})/M_s\).
Keep this improving factor in
\eqref{eq:hall:uniform-forcing} and add the high-\(H\) tail.
Optimizing gives, for \(a>0\), a bound on the normalized magnitude of
the \emph{signed} endpoint contribution of
\begin{equation}
                       C(\epsilon+r^{-1})^{1/(a+1)} .
 \label{eq:hall:quantum-end}
\end{equation}
Here the normalization is \(r^{p_G}(1+\log r)^d\). The bound
includes a finite union of endpoint bands, possibly shrinking with
the level. It is distinct from the positive classical exhaustion
\eqref{eq:hall:classical-end}.

At fixed forced \(H\), choose a dormant bipartition so that a given
dormant endpoint is unsummed. Its folded harmonic strips
\(n<r^\zeta\) and \(n>r^{1-\zeta}\) cost normalized
\(O(\zeta+1/\log r)\). For intersections choose one such restriction
as unsummed and use the signed prefix estimates for the other
ones. If an originally reflected dormant site lies in the remaining
mesoscopic rectangle, its untouched \(1-\cos\) supplies
\(O(r^{-2\zeta})\) in the same orientation. Thus only the original
negative mesoscopic rectangle remains. In a strict violation of one
dormant Hall inequality by logarithmic gap \(\gamma>0\), the factor
\(n_j/M_{b_j}\) loses \(O(r^{-\gamma})\); thin comparison strips cost
their log width. On a strict allowed chamber each dormant matched
anchor dominates its incident labels by a power, including the
bounded forced ones. The signed tail estimate removes matched
products \(w_j>K\) at normalized \(O(1/K)\) cost. The volume of
\(w_j<1/K\) is \(O(1/K)+o(1)\) by the odd mesh on separated log
cells. All of these cutoffs have now been justified in the
original exact signed diagram, prior to any approximation of an
attaching matrix.

In a retained compact product cell each \(T_*\)--\(I\) attaching
monodromy tends to the identity in the unitary ribbon category.
For fixed forced color \(n\le H\) and dormant deficit \(m=o(r)\),
the balancing phase on a fused spin \(j+s\), \(|s|\le(n-1)/2\),
differs from one by \(O((Hm+H^2)/r)\). This bounds the operator
norm because fusion channels are orthogonal. Replacing these
finitely many monodromies therefore changes the normalized compact
cell by \(o(1)\): the normalized quantum traces and the bounded
fixed packets do not amplify the error, and the absolute
compact-annulus radial count was already bounded. There is no
absolute summation of this perturbation over unbounded products.
With the chosen exterior routing the resulting old sublink is
split and undecorated. The old macro bracket remains exactly the
ordinary normalized Hopf bracket at its original values. The
dormant calculation is now on the disconnected heavy forest with
all original fixed auxiliary contractions, so the actual
coherent-state kernel of Theorem~\ref{thm:ang:compact} retains their
joint \(P_W\). This proves the fixed-prefix assertion in
Theorem~\ref{thm:hall:overlap}, including connecting words.

For each fixed old odd color \(n\) the normalized bracket is
\[
 \frac{\sin(\pi/r)}{\sin(\pi n/r)}
                      \frac{\sin(\pi n x)}{\sin(\pi x)} .
\]
It is uniformly bounded for the fixed \(n\) on the full macro
interval, with removable values at both endpoints. Thus its Riemann
limit with the actual cap factor is the ordinary normalized
character integral. The limit as \(H\to\infty\) is
\eqref{eq:hall:weight} by \eqref{eq:hall:characters} on compact
macro windows. On the full box the same identification follows
from \eqref{eq:hall:uniform-forcing},
\eqref{eq:hall:quantum-end}, and the positive exhaustion
\eqref{eq:hall:classical-end}. This removes the old macro and
fixed hub gates without referring to a limiting normalized active
angle at a degenerate cap relator.

Finally use the previously controlled signed dormant tails to
let the compact products exhaust the positive half-lines in the
strict Hall chamber. In the bounded-product joint kernel,
independent cap sign flips realize any assignment of signs to
the edge dots of a forest and leave \(P_W\) invariant. Thus the
matched Fourier integrations may be taken with real cosines.
For one matched edge the Abel kernel
\[
       \int_0^\infty e^{-\epsilon w}\cos(\pi wt)\,dw
          =\frac{\epsilon}{\epsilon^2+\pi^2t^2}
\]
tends to \(\delta_0(t)\) with mass one on \(\mathbb R\). The signed
tail bound identifies this auxiliary Abel value with the
ungated physical leading integral. Each bare dot probability
\(dt/2\) contributes \(1/2\), canceling the \(2^d\) in
\eqref{eq:hall:overlap}. Conditional Haar at the matched zero
dots gives precisely the joint frames specified in the theorem.
Exhausting the compact chambers and the old forced tuples yields
\eqref{eq:hall:main}. This is also the classical absolute
normalization in \eqref{eq:hall:normalization}, obtained by a
different computation.
\end{proof}

The genuine joint spin-one words of Section~\ref{sec:ang:main}
therefore determine the normalized dormant law on its
unprojected frame quotient, even across components in the
compatible routing. For example, a single edge-band spin-one
character \(4t^2-1\) has normalized leading mean \(-1\) on a
matched edge and \(1/3\) on an unmatched dormant edge; products
of distinct non-edge observables must still be taken in the shared
frame law. In contrast, positive physical product gates
\(\psi\) supported in \(1<w<2\) on a matched edge can give a
negative vacuum overlap coefficient in
\eqref{eq:hall:overlap}, from the negative sinc on that interval.
The positive frame measure after signed integration is therefore
not a positive joint measure of unintegrated product and angle
data. The theorem does not include arbitrary active angular
Wilson insertions, a route that truly links the separated forced
sublink, or growing packet colors. Nor does it assign lower-order
coefficients of a globally ungated integral to all other
character strata.

\subsection{A deficient five-vertex validation}
\label{sec:hall:broom}

The first mixed diagram with a nontrivial dormant matching is
the five-vertex tree with edges
\(s_1c,s_2c,cd,de\). Its intrinsic sets are
\(S_*=\{s_1,s_2\}, T_*=\{c\}, I=\{d,e\}\).
Here \(p_G=4\) and \(c_G=2/\pi^4\). At fixed compact old macro radii
the measured dormant regular negative branch has half-torsion
\[
    32\pi^4\,dM_*\,
       \frac{d\Pi(u_d)d\Pi(u_e)}{|u_d\cdot u_e|}.
\]
At the perpendicular overlap the replacement by \(\pi^2\,dy_d\,dy_e\)
and the calibrated cubic integration give
\((2/\pi^4)r^4\log r\,dM_*\).
This is also a direct finite check without the packet theorem.
The two compact old macro transfers decay rapidly in the odd hub
index by ordinary sine summation. At fixed hub index \(n\), the
remaining terminal transfer is
\((1+O(1/m))/m\) at its adjacent dormant label \(m\) in negative
trace windows equal to one at their endpoints. Its middle sinc
changes the harmonic sum by \(O(1+\log(1+n))\), so the coefficient
is the odd harmonic half-log times the old forced character
sum. The latter is exactly \(\int dM_*\) with the chosen macro
tests. This recovers the absolute coefficient and a remainder
\(O(r^4)\).

As a normalization check across a true positive end, let
\(h_d=-1,h_e=+1\). For each fixed positive \(e\) root the dormant
adjacency on the line and full Lie algebra has nonzero eigenvalues
\(\pm1\). Its kernel is the actual plane
\(v=\log(h_e)/\pi\perp u_d\); by
\eqref{eq:hall:torsion} its boundary density is
\(2\pi^5dM_*\,d\Pi(u_d)\,d^2v\).
There are two positive root choices. Writing \(y=1-|v|\), the
continued activated cover has exact radial stack density
\(16\pi^5\cot(\pi y/2)\,dy\,dM_*\), whose endpoint agrees with
the plane density after polar integration. With a smooth literal
upper-end mark \(\eta_e(y)\) and the other dormant and hub
negative windows as above, the exact quantum leading term is
\[
      \frac{r^4}{\pi^3}
           \left(\int\eta\,dM_*\right)
           \left(\int_0^1\eta_e(y)\cot(\pi y/2)\,dy\right)
      +O(r^3).
\]
Indeed the \(e\)-transfer at upper \(y\) is the sine transform of the
odd-reflected smooth cap profile and is rapidly summable in
the remaining lower index. The attaching sinc can be replaced by
one at \(O(r^{-1})\) cost in the normalized sum. The odd sine
series gives the e integral and the same forced character
factor as before. There is no new logarithmic residue.

If the other dormant site is positive instead, its opening changes
the hub product to the actual three-input ordered product. For the
activated cover write
\[
 dM'=\left(\prod_{s=s_1,s_2,d}2\cos^2(\pi x_s/2)\,dx_s\right)
          L_c(x_{s_1},x_{s_2},x_d)\,L_e(x_d),
 \qquad L_e(y)=\frac{\pi}{4\sin(\pi y)} .
\]
For a literal \(d\) upper-end mark \(\eta_d\) and \(e,c\) negative
windows the corresponding coefficient is
\[
                4\pi^{-4}r^4\int\eta\,\eta_d\,dM'+O(r^3).
\]
One can verify the exact coefficient directly: fold the upper \(d\)
label \(m\) to \(r-m\) in the terminal transfer, whose error is
\(O((r-m)^{-1})\). For odd hub \(n\) the middle sine is
\(O(n(r-m)/r)\), so this error is summable at \(O(r^{-1})\) in the
normalized sum. Its odd-mesh integral supplies the third normalized
character at \(c\) with the factor \(2\cos^2(\pi y/2)L_e(y)\).
The two old compact macro profiles make the resulting triple
character series summable all the way to the positive \(d\) endpoint.

If both dormant sites are positive they form a separate
four-dimensional smooth chart: their full projected adjacency
is invertible and there is no nearby noncentral dormant site.
Equation~\eqref{eq:hall:torsion} gives total stack
\((\pi^2/2)dM_*\) after the four positive roots. Its action is
\(1/2\) modulo one. With both positive windows equal to one at
their endpoints the exact coefficient is
\((-1)^r r^2\int\eta\,dM_*/(8\pi^2)+O(r)\).
For the finite check fold both dormant labels and write
\(m=r-i\) on \(d\). The terminal sine transfer has the double
reflection factor \((-1)^r\) and the folded transform of
\[
                 H(z)=\frac{\pi}{2}\eta_e(1-z)\tan(\pi z/2).
\]
The middle sinc supplies \(i/r+O(r^{-2})\) at fixed odd hub \(n\)
and fixed \(i\), and the remaining dormant profile supplies
\(\pi^2/(4r)+O(r^{-2})\). The fixed-parity sine derivative identity
\[
 \sum_{\substack{i>0\\ i\equiv r-1\ (2)}} i\,\frac2\pi
      \int_0^1 H(z)\sin(\pi i z)\,dz
            =\frac{H'(0)}{2\pi^2}=\frac18
\]
follows by reflecting the sine series and using \(H'(1)=0\).
Smooth odd continuation at the reflected end makes the weighted
transforms summable uniformly, so the same estimates give the
indicated remainder. This separates the two-positive chart from
the one-positive smooth endings of leading regular covers; it
is not an additional Hall logarithmic residue.
